International Competition in Mathematics for
             Universtiy Students
                      in
              Plovdiv, Bulgaria
                    1996
                                                                             1

PROBLEMS AND SOLUTIONS

   First day — August 2, 1996


   Problem 1. (10 points)
   Let for j = 0, . . . , n, aj = a0 + jd, where a0 , d are fixed real numbers.
Put                                                      
                             a0 a1      a2     . . . an
                            a
                            1 a0       a1     . . . an−1 
                                                          
                   A =  a2 a1          a0     . . . an−2  .
                                                         
                            ........................... 
                                                         

                             an an−1 an−2 . . . a0

Calculate det(A), where det(A) denotes the determinant of A.
   Solution. Adding the first column of A to the last column we get that
                                                               
                                     a0 a1      a2    ... 1
                                    a
                                    1 a0       a1    ... 1 
           det(A) = (a0 + an ) det  a2 a1      a0    ... 1 .
                                                            
                                    ....................... 
                                                            

                                     an an−1 an−2 . . . 1

Subtracting the n-th row of the above matrix from the (n+1)-st one, (n−1)-
st from n-th, . . . , first from second we obtain that
                                                             
                                       a0 a1 a2 . . . 1
                                      d −d −d . . . 0 
                                                           
             det(A) = (a0 + an ) det  d     d −d . . . 0  .
                                                           
                                      .................... 
                                                           

                                        d    d    d ... 0

Hence,
                                                                 
                                        d −d −d . . . −d
                                       d   d −d . . . −d 
                                                            
                      n
         det(A) = (−1) (a0 + an ) det  d   d    d . . . −d  .
                                                            
                                       .................... 
                                                            

                                        d   d    d ...     d
2

Adding the last row of the above matrix to the other rows we have
                                                              
                              2d 0 0 . . . 0
                             2d 2d     0 ... 0 
                                                 
             n
det(A) = (−1) (a0 +an ) det  2d 2d 2d . . . 0  = (−1)n (a0 +an )2n−1 dn .
                                                 
                             ................... 
                                                 

                               d d d ... d

    Problem 2. (10 points)
    Evaluate the definite integral
                             Z π
                                    sin nx
                                       x
                                               dx,
                              −π (1 + 2 )sin x

where n is a natural number.
    Solution. We have
                     Z π
                           sin nx
             In =             x
                                     dx
                    −π (1 + 2 )sin x
                   Z π                    Z 0
                          sin nx                 sin nx
                 =            x
                                    dx  +           x
                                                            dx.
                    0 (1 + 2 )sin x        −π (1 + 2 )sin x

In the second integral we make the change of variable x = −x and obtain
                     Z π                            Z π
                           sin nx                             sin nx
             In =             x
                                      dx +                                dx
                    0 (1 + 2 )sin x                   0   (1 + 2−x )sin x
                       (1 + 2x )sin nx
                   Z π
                 =                     dx
                    0   (1 + 2x )sin x
                   Z π
                       sin nx
                 =            dx.
                    0   sin x
For n ≥ 2 we have
                                   Z π
                                              sin nx − sin (n − 2)x
                 In − In−2 =                                        dx
                                      0               sin x
                                      Z π
                             = 2               cos (n − 1)xdx = 0.
                                          0

The answer                        (
                                      0        if n is even,
                           In =
                                      π        if n is odd
                                                                            3

follows from the above formula and I 0 = 0, I1 = π.
   Problem 3. (15 points)
   The linear operator A on the vector space V is called an involution if
A2 = E where E is the identity operator on V . Let dim V = n < ∞.
   (i) Prove that for every involution A on V there exists a basis of V
consisting of eigenvectors of A.
   (ii) Find the maximal number of distinct pairwise commuting involutions
on V .
      Solution.
                 1
      (i) Let B = (A + E). Then
                 2
                  1 2              1            1
           B2 =     (A + 2AE + E) = (2AE + 2E) = (A + E) = B.
                  4                4            2
Hence B is a projection. Thus there exists a basis of eigenvectors for B, and
the matrix of B in this basis is of the form diag(1, . . . , 1, 0, . . . , 0).
    Since A = 2B − E the eigenvalues of A are ±1 only.
    (ii) Let {Ai : i ∈ I} be a set of commuting diagonalizable operators
on V , and let A1 be one of these operators. Choose an eigenvalue λ of A 1
and denote Vλ = {v ∈ V : A1 v = λv}. Then Vλ is a subspace of V , and
since A1 Ai = Ai A1 for each i ∈ I we obtain that Vλ is invariant under each
Ai . If Vλ = V then A1 is either E or −E, and we can start with another
operator Ai . If Vλ 6= V we proceed by induction on dim V in order to find
a common eigenvector for all Ai . Therefore {Ai : i ∈ I} are simultaneously
diagonalizable.
    If they are involutions then |I| ≤ 2n since the diagonal entries may equal
1 or -1 only.
      Problem 4. (15 points)
                       1 n−1
                         X
      Let a1 = 1, an =       ak an−k for n ≥ 2. Show that
                       n k=1
      (i) lim sup |an |1/n < 2−1/2 ;
           n→∞
      (ii) lim sup |an |1/n ≥ 2/3.
           n→∞
      Solution.
      (i) We show by induction that

(∗)                              an ≤ q n   for   n ≥ 3,
4

                                                               1         1
where q = 0.7 and use that 0.7 < 2−1/2 . One has a1 = 1, a2 = , a3 = ,
                                                               2         3
     11
a4 =    . Therefore (∗) is true for n = 3 and n = 4. Assume (∗) is true for
     48
n ≤ N − 1 for some N ≥ 5. Then

       2        1        1 NX
                            −3
                                         2        1         N −5 N
aN =     aN −1 + aN −2 +       ak aN −k ≤ q N −1 + q N −2 +     q ≤ qN
       N        N        N k=3           N        N          N

           2   1
because      + 2 ≤ 5.
           q q
     (ii) We show by induction that
                                  an ≥ q n       for       n ≥ 2,
                                          2
               2               1             2
where q =        . One has a2 = >                      = q 2 . Going by induction we have
               3               2             3
for N ≥ 3

              2        1 NX
                          −2
                                       2         N −3 N
          aN = aN −1 +       ak aN −k ≥ q N −1 +     q = qN
              N        N k=2           N          N

         2
because    = 3.
         q
   Problem 5. (25 points)
   (i) Let a, b be real numbers such that b ≤ 0 and 1 + ax + bx 2 ≥ 0 for
every x in [0, 1]. Prove that
                                               1
                                                             
                            Z 1
                                                 if a < 0,       −
                                                             
                                             2 n
                lim n   (1 + ax + bx ) dx =    a
               n→+∞   0                      +∞ if a ≥ 0.
    (ii) Let f : [0, 1] → [0, ∞) be a function with a continuous second
derivative and let f 00 (x) ≤ 0 for every x in [0, 1]. Suppose that L =
        Z 1
lim n         (f (x))n dx exists and 0 < L < +∞. Prove that f 0 has a con-
n→∞       0
stant sign and min |f 0 (x)| = L−1 .
                  x∈[0,1]
     Solution. (i) With a linear change of the variable (i) is equivalent to:
     (i0 ) Let a, b, A be real numbers such that b ≤ 0, A > 0 and 1+ax+bx 2 > 0
                                                   Z A
for every x in [0, A]. Denote In = n                       (1 + ax + bx2 )n dx. Prove that
                                                       0
                 1
    lim In = −     when a < 0 and lim In = +∞ when a ≥ 0.
n→+∞             a               n→+∞
                                                                                                              5

   Let a < 0. Set f (x) = eax − (1 + ax + bx2 ). Using that f (0) = f 0 (0) = 0
and f 00 (x) = a2 eax − 2b we get for x > 0 that

                                     0 < eax − (1 + ax + bx2 ) < cx2

             a2
where c =       − b. Using the mean value theorem we get
             2

                            0 < eanx − (1 + ax + bx2 )n < cx2 nea(n−1)x .

Therefore
            Z A                              Z A                                      Z A
                            anx                                 2 n               2
      0<n               e         dx − n           (1 + ax + bx ) dx < cn                   x2 ea(n−1)x dx.
                0                             0                                        0

Using that
                                                           eanA − 1
                                        Z A
                                                                         1
                                    n         eanx dx =             −→ −
                                         0                    a     n→∞ a

and                         Z A
                                                                 1
                                                                            Z ∞
                                   x2 ea(n−1)x dx <                               t2 e−t dt
                             0                            |a|3 (n − 1)3      0

we get (i0 ) in the case a < 0.
   Let a ≥ 0. Then for n > max{A−2 , −b} − 1 we have

                Z A                                            Z √1
                                               2 n                    n+1
            n           (1 + ax + bx ) dx > n                               (1 + bx2 )n dx
                    0                                           0
                                                                                                 n
                                                                1          b
                                                                                 
                                                        > n· √       · 1+
                                                               n+1        n+1
                                                             n     b
                                                        > √       e −→ ∞.
                                                           n + 1 n→∞

(i) is proved.
                                     Z 1
   (ii) Denote In = n                        (f (x))n dx and M = max f (x).
                                        0                              x∈[0,1]
   For M < 1 we have In ≤ nM n −→ 0, a contradiction.
                                                       n→∞
    If M > 1 since f is continuous there exists an interval I ⊂ [0, 1] with
|I| > 0 such that f (x) > 1 for every x ∈ I. Then I n ≥ n|I| −→ +∞,
                                                                                                      n→∞
a contradiction. Hence M = 1. Now we prove that f 0 has a constant
sign. Assume the opposite. Then f 0 (x0 ) = 0 for some x ∈ (0, 1). Then
6

                                                                                           h2 00
f (x0 ) = M = 1 because f 00 ≤ 0. For x0 + h in [0, 1], f (x0 + h) = 1 +                     f (ξ),
                                                                                    2
                                                                                           2
ξ ∈ (x0 , x0 + h). Let m = min f 00 (x). So, f (x0 + h) ≥ 1 + h2 m.
                                    x∈[0,1]
                                           δ2
     Let δ > 0 be such that 1 +               m > 0 and x0 + δ < 1. Then
                                           2
                        Z x0 +δ                     Z δ                   n
                                                                  m 2
               In ≥ n             (f (x))n dx ≥ n          1+       h           dh −→ ∞
                            x0                      0             2                n→∞


in view of (i0 ) – a contradiction. Hence f is monotone and M = f (0) or
M = f (1).
    Let M = f (0) = 1. For h in [0, 1]
                                                                           m 2
                            1 + hf 0 (0) ≥ f (h) ≥ 1 + hf 0 (0) +            h ,
                                                                           2

where f 0 (0) 6= 0, because otherwise we get a contradiction as above. Since
f (0) = M the function f is decreasing and hence f 0 (0) < 0. Let 0 < A < 1
                             m
be such that 1 + Af 0 (0) + A2 > 0. Then
                             2
     Z A                             Z A                     Z A                              n
                                                                                         m 2
 n         (1 + hf 0 (0))n dh ≥ n          (f (x))n dx ≥ n             1 + hf 0 (0) +      h        dh.
      0                               0                       0                          2

                                                                              1
From (i0 ) the first and the third integral tend to −                              as n → ∞, hence
                                                                           f 0 (0)
so does theZ second.
                 1                                                                                 1
     Also n          (f (x))n dx ≤ n(f (A))n −→ 0 (f (A) < 1). We get L = −
                A                                n→∞                                            f 0 (0)
in this case.
                                                                  1
    If M = f (1) we get in a similar way L =                           .
                                                             f 0 (1)
     Problem 6. (25 points)
     Upper content of a subset E of the plane R 2 is defined as
                                                ( n                   )
                                                 X
                                  C(E) = inf            diam(Ei )
                                                  i=1


where inf is taken over all finite families of sets E 1 , . . . , En , n ∈ N, in R2
                n
such that E ⊂ ∪ Ei .
                      i=1
                                                                                    7

     Lower content of E is defined as

           K(E) = sup {lenght(L)       : L is a closed line segment
                                          onto which E can be contracted} .

Show that
    (a) C(L) = lenght(L) if L is a closed line segment;
    (b) C(E) ≥ K(E);
    (c) the equality in (b) needs not hold even if E is compact.
    Hint. If E = T ∪ T 0 where T is the triangle with vertices (−2, 2), (2, 2)
and (0, 4), and T 0 is its reflexion about the x-axis, then C(E) = 8 > K(E).
    Remarks: All distances used in this problem are Euclidian. Diameter
of a set E is diam(E) = sup{dist(x, y) : x, y ∈ E}. Contraction of a set E
to a set F is a mapping f : E 7→ F such that dist(f (x), f (y)) ≤ dist(x, y) for
all x, y ∈ E. A set E can be contracted onto a set F if there is a contraction
f of E to F which is onto, i.e., such that f (E) = F . Triangle is defined as
the union of the three segments joining its vertices, i.e., it does not contain
the interior.
     Solution.
     (a) The choice E1 = L gives C(L) ≤ lenght(L). If E ⊂ ∪ni=1 Ei then
n
X
    diam(Ei ) ≥ lenght(L): By induction, n=1 obvious, and assuming that
i=1
En+1 contains the end point a of L, define the segment L ε = {x ∈ L :
                                                             n+1
                                                             X
dist(x, a) ≥ diam(En+1 )+ε} and use induction assumption to get              diam(Ei ) ≥
                                                                       i=1
lenght(Lε ) + diam(En+1 ) ≥ lenght(L) − ε; but ε > 0 is arbitrary.
    (b) If f is a contraction of E onto L and E ⊂ ∪ nn=1 Ei , then L ⊂ ∪ni=1 f (Ei )
                      n
                      X                       n
                                              X
and lenght(L) ≤             diam(f (Ei )) ≤         diam(Ei ).
                      i=1                     i=1
   (c1) Let E = T ∪ T 0 where T is the triangle with vertices (−2, 2), (2, 2)
                                                                       n
and (0, 4), and T 0 is its reflexion about the x-axis. Suppose E ⊂ ∪ Ei .
                                                                              i=1
If no set among Ei meets both T and T 0 , then Ei may be partitioned into
covers of segments [(−2, 2), (2, 2)] and [(−2, −2), (2, −2)], both of length 4,
     n
     X
so         diam(Ei ) ≥ 8. If at least one set among Ei , say Ek , meets both T and
     i=1
T 0 , choose a ∈ Ek ∩ T and b ∈ Ek ∩ T 0 and note that the sets Ei0 = Ei for
i 6= k, Ek0 = Ek ∪ [a, b] cover T ∪ T 0 ∪ [a, b], which is a set of upper content
8

at least 8, since its orthogonal projection onto y-axis is a segment of length
                                          n
                                          X
8. Since diam(Ej ) = diam(Ej0 ), we get         diam(Ei ) ≥ 8.
                                          i=1
    (c2) Let f be a contraction of E onto L = [a 0 , b0 ]. Choose a = (a1 , a2 ),
b = (b1 , b2 ) ∈ E such that f (a) = a0 and f (b) = b0 . Since lenght(L) =
dist(a0 , b0 ) ≤ dist(a, b) and since the triangles have diameter only 4, we may
assume that a ∈ T and b ∈ T 0 . Observe that if a2 ≤ 3 then a lies on one of
the segments joining some of the points (−2, 2), (2, 2), (−1, 3), (1, 3); since
all these
      √ points have distances from vertices, and       √ so from points, of T 2 at
most 50, we get that lenght(L) ≤ dist(a, b) ≤ 50. Similarly if b2 ≥ −3.
Finally, if a2 > 3 and b2 < −3, we√note that every vertex, and so every point
of T is in the distance at most 10 for√a and every vertex, and so every
point, of T 0 is in the distance at most 10 of b. Since f is a contraction,
                                                                        √
the image of T lies in a segment containing a 0 of length at most     √   10 and
the image of T 0 lies in a segment containing b0 of length at√most √10. Since
the union√of these two images is L, we get lenght(L) ≤ 2 10 ≤ 50. Thus
K(E) ≤ 50 < 8.

    Second day — August 3, 1996

    Problem 1. (10 points)
    Prove that if f : [0, 1] → [0, 1] is a continuous function, then the sequence
of iterates xn+1 = f (xn ) converges if and only if

                              lim (xn+1 − xn ) = 0.
                             n→∞


     Solution. The “only if” part is obvious. Now suppose that lim (xn+1
                                                                      n→∞
−xn ) = 0 and the sequence {xn } does not converge. Then there are two
cluster points K < L. There must be points from the interval (K, L) in the
                                                                   |f (x) − x|
sequence. There is an x ∈ (K, L) such that f (x) 6= x. Put ε =                 >
                                                                        2
0. Then from the continuity of the function f we get that for some δ > 0 for
all y ∈ (x−δ, x+δ) it is |f (y)−y| > ε. On the other hand for n large enough
it is |xn+1 − xn | < 2δ and |f (xn ) − xn | = |xn+1 − xn | < ε. So the sequence
cannot come into the interval (x − δ, x + δ), but also cannot jump over this
interval. Then all cluster points have to be at most x − δ (a contradiction
with L being a cluster point), or at least x + δ (a contradiction with K being
a cluster point).
                                                                           9

      Problem 2. (10 points)
                                                     et + e−t
      Let θ be a positive real number and let cosh t =        denote the
                                                         2
hyperbolic cosine. Show that if k ∈ N and both cosh kθ and cosh (k + 1)θ
are rational, then so is cosh θ.
      Solution. First we show that

(1)        If cosh t is rational and m ∈ N, then cosh mt is rational.

Since cosh 0.t = cosh 0 = 1 ∈ Q and cosh 1.t = cosh t ∈ Q, (1) follows
inductively from

               cosh (m + 1)t = 2cosh t.cosh mt − cosh (m − 1)t.

    The statement of the problem is obvious for k = 1, so we consider k ≥ 2.
For any m we have
(2)
   cosh θ = cosh ((m + 1)θ − mθ) =
          = cosh (m + 1)θ.cosh mθ − sinh
                                     p
                                         (m + 1)θ.sinh mθ √
          = cosh (m + 1)θ.cosh mθ − cosh 2 (m + 1)θ − 1. cosh 2 mθ − 1

Set cosh kθ = a, cosh (k + 1)θ = b, a, b ∈ Q. Then (2) with m = k gives
                                        p        p
                        cosh θ = ab −       a2 − 1 b2 − 1

and then
                 (a2 − 1)(b2 − 1) = (ab − cosh θ)2
(3)
                                  = a2 b2 − 2abcosh θ + cosh 2 θ.

   Set cosh (k 2 − 1)θ = A, cosh k 2 θ = B. From (1) with m = k − 1 and
t = (k + 1)θ we have A ∈ Q. From (1) with m = k and t = kθ we have
B ∈ Q. Moreover k 2 − 1 > k implies A > a and B > b. Thus AB > ab.
From (2) with m = k 2 − 1 we have

               (A2 − 1)(B 2 − 1) = (AB − cosh θ)2
(4)
                                 = A2 B 2 − 2ABcosh θ + cosh 2 θ.

    So after we cancel the cosh 2 θ from (3) and (4) we have a non-trivial
linear equation in cosh θ with rational coefficients.
10

     Problem 3. (15 points)
     Let G be the subgroup of GL2 (R), generated by A and B, where
                         "               #               "               #
                             2       0                       1       1
                    A=                       ,   B=                          .
                             0       1                       0       1
                                                         !
                                   a11 a12
Let H consist of those matrices               in G for which a11 =a22 =1.
                                   a21 a22
   (a) Show that H is an abelian subgroup of G.
   (b) Show that H is not finitely generated.
    Remarks. GL2 (R) denotes, as usual, the group (under matrix multipli-
cation) of all 2 × 2 invertible matrices with real entries (elements). Abelian
means commutative. A group is finitely generated if there are a finite number
of elements of the group such that every other element of the group can be
obtained from these elements using the group operation.
     Solution.
     (a) All of the matrices in G are of the form
                                     "               #
                                         ∗       ∗
                                                         .
                                         0       ∗

So all of the matrices in H are of the form
                                             "               #
                                                 1       x
                             M (x) =                             ,
                                                 0       1

so they commute. Since M (x)−1 = M (−x), H is a subgroup of G.
    (b) A generator of H can only be of the form M (x), where x is a binary
                     p
rational, i.e., x = n with integer p and non-negative integer n. In H it
                    2
holds

                             M (x)M (y) = M (x + y)
                          M (x)M (y)−1 = M (x − y).

                                1
                                        
The matrices of the form M         are in H for all n ∈ N. With only finite
                               2n
number of generators all of them cannot be achieved.
                                                                          11

     Problem 4. (20 points)
     Let B be a bounded closed convex symmetric (with respect to the origin)
set in R2 with boundary the curve Γ. Let B have the property that the
ellipse of maximal area contained in B is the disc D of radius 1 centered at
the origin with boundary the circle C. Prove that A ∩ Γ 6= Ø for any arc A
                      π
of C of length l(A) ≥ .
                       2
     Solution. Assume the contrary – there is an arc A ⊂ C with length
         π
l(A) = such that A ⊂ B\Γ. Without loss of generality we may assume that
         2                  √     √              √       √
the ends of A are M = (1/ 2, 1/ 2), N = (1/ 2, −1/ 2). A is compact
and Γ is closed. From A ∩ Γ = Ø we get δ > 0 such that dist(x, y) > δ for
every x ∈ A, y ∈ Γ.
                                                             x2      y2
Given ε > 0 with Eε we denote the ellipse with boundary:          2
                                                                    + 2 = 1,
                                                          (1 + ε)    b
such that M, N ∈ Eε . Since M ∈ Eε we get
                                     (1 + ε)2
                            b2 =                 .
                                   2(1 + ε)2 − 1
Then we have
                                  (1 + ε)2
                 area Eε = π p                 > π = area D.
                                 2(1 + ε)2 − 1
In view of the hypotheses, Eε \ B 6= Ø for every ε > 0. Let S = {(x, y) ∈
R2 : |x| > |y|}. ¿From Eε \ S ⊂ D ⊂ B it follows that Eε \ B ⊂ S. Taking
ε < δ we get that
                   Ø 6= Eε \ B ⊂ Eε ∩ S ⊂ D1+ε ∩ S ⊂ B
– a contradiction (we use the notation D t = {(x, y) ∈ R2 : x2 + y 2 ≤ t2 }).
    Remark. The ellipse with maximal area is well known as John’s ellipse.
Any coincidence with the President of the Jury is accidental.
   Problem 5. (20 points)
   (i) Prove that
                                 ∞
                                 X        nx           1
                           lim                     =     .
                         x→+∞
                                 n=1
                                       (n2 + x)2       2
(ii) Prove that there is a positive constant c such that for every x ∈ [1, ∞)
we have
                            ∞
                           X        nx      1     c
                                  2     2
                                          −    ≤ .
                           n=1
                               (n + x)      2     x
12

      Solution.
                              t             1
      (i) Set f (t) =           2  2
                                     , h = √ . Then
                          (1 + t )           x
                   ∞                 ∞
                          nx                         ∞          1
                   X                X                            Z
                        2     2
                                = h     f (nh) −→      f (t)dt = .
                  n=1
                      (n + x)       n=1
                                               h→0 0            2

                                                ∞
                                                X
      The convergence holds since h                   f (nh) is a Riemann sum of the inte-
       Z ∞                                      n=1

gral         f (t)dt. There are no problems with the infinite domain because
        0
                                                                      ∞
                                                                      X                   Z ∞
f is integrable and f ↓ 0 for x → ∞ (thus h                                  f (nh) ≥                  f (t)dt ≥
                                                                      n=N                     nN
      ∞
      X
h            f (nh)).
    n=N +1
      (ii) We have
          ∞                   ∞                               Z nh+ h            !        Z h
          X      nx      1   X                                          2                         2
                       −   =     hf (nh) −                                  f (t)dt −                 f (t)dt
          n=1
              (n2 + x)2 2    n=1                               nh− h2                         0
(1)                                  ∞                      Z nh+ h                  Z h
                                     X                            2                       2
                                 ≤            hf (nh) −                 f (t)dt +             f (t)dt
                                     n=1                      nh− h
                                                                  2
                                                                                      0

      Using twice integration by parts one has
                        Z a+b                  Z b
                                  1
(2)       2bg(a) −     g(t)dt = −                    (b − t)2 (g 00 (a + t) + g 00 (a − t))dt
                   a−b            2             0

for every g ∈ C 2 [a − b, a + b]. Using f (0) = 0, f ∈ C 2 [0, h/2] one gets
                                   Z h/2
(3)                                           f (t)dt = O(h2 ).
                                     0

      From (1), (2) and (3) we get
        ∞                   ∞                 Z nh+ h
                nx      1                              2
                               h2                          |f 00 (t)|dt + O(h2 ) =
        X                  X
               2 + x)2
                       − ≤
        n=1
            (n          2  n=1                 nh− h
                                                   2
                                         Z ∞
                                = h2           |f 00 (t)|dt + O(h2 ) = O(h2 ) = O(x−1 ).
                                          h
                                          2
                                                                                                   13

      Problem 6. (Carleman’s inequality) (25 points)
      (i) Prove that for every sequence {a n }∞
                                              n=1 , such that an > 0, n = 1, 2, . . .
       ∞
       X
and         an < ∞, we have
      n=1

                               ∞                                      ∞
                                     (a1 a2 · · · an )1/n < e
                               X                                      X
                                                                            an ,
                               n=1                                    n=1

where e is the natural log base.
   (ii) Prove that for every ε > 0 there exists a sequence {a n }∞
                                                                 n=1 , such that
                             ∞
                             X
an > 0, n = 1, 2, . . .,            an < ∞ and
                             n=1

                            ∞                                             ∞
                                  (a1 a2 · · · an )1/n > (e − ε)
                            X                                             X
                                                                                an .
                            n=1                                           n=1

      Solution.
      (i) Put for n ∈ N

(1)                                       cn = (n + 1)n /nn−1 .

Observe that c1 c2 · · · cn = (n + 1)n . Hence, for n ∈ N,

                 (a1 a2 · · · an )1/n = (a1 c1 a2 c2 · · · an cn )1/n /(n + 1)
                                            ≤ (a1 c1 + · · · + an cn )/n(n + 1).

Consequently,
               ∞                                ∞                ∞
                                                                                               !
                     (a1 a2 · · · an )1/n ≤
               X                                X                X
(2)                                                  an cn             (m(m + 1))−1 .
               n=1                             n=1            m=n

Since              ∞                               ∞ 
                                                       1               1
                   X                               X                               
                         (m(m + 1))−1 =                           −                    = 1/n
                 m=n                              m=n        m        m+1
we have                                                           !
                      ∞
                      X               ∞
                                      X                                   ∞
                                                                          X
                                                             −1
                            an cn           (m(m + 1))                =            an cn /n
                      n=1            m=n                                  n=1
                                    ∞                                  ∞
                                          an ((n + 1)/n)n < e
                                    X                                  X
                              =                                               an
                                    n=1                                n=1
14

(by (1)). Combining the last inequality with (2) we get the result.
   (ii) Set an = nn−1 (n + 1)−n for n = 1, 2, . . . , N and an = 2−n for n > N ,
where N will be chosen later. Then
                                                          1
(3)                               (a1 · · · an )1/n =
                                                         n+1
for n ≤ N . Let K = K(ε) be such that
                                     n
                               n+1                   ε
                           
(4)                                        >e−         for n > K.
                                n                    2
Choose N from the condition
                 K             ∞                                  N
                 X             X                       ε         X    1
(5)                    an +          2−n ≤                              ,
                 n=1           n=1
                                                (2e − ε)(e − ε) n=K+1 n

which is always possible because the harmonic series diverges. Using (3), (4)
and (5) we have
         ∞           K             ∞                   N                  n
                                                            1        n
         X           X             X                   X        
                                               −n
              an =         an +            2        +                          <
        n=1          n=1          n=N +1              n=K+1
                                                            n       n+1
                                    N                               −1    N
                         ε              1      ε                                1
                                   X                                      X
                <                         + e−                                    =
                  (2e − ε)(e − ε) n=K+1 n      2                          n=K+1
                                                                                n
                             N              ∞
                       1         1     1 X
                                               (a1 · · · an )1/n .
                            X
                =                  ≤
                     e − ε n=K+1 n   e − ε n=1
